% =============================================================
%  D3 — Scalable Training & Distributed Computing
%  MLOps Project — ResNet50 / FER2013
%  Paste this section into your main Overleaf document.
%  Requires: booktabs, amsmath, graphicx, listings, xcolor
% =============================================================

\section{D3 — Scalable Training and Distributed Computing}

\subsection{Introduction}

Training deep neural networks efficiently at scale requires understanding both the
theoretical limits of parallelism and the practical trade-offs of different hardware
and software strategies. This section documents the distributed training experiments
conducted on the FER2013 facial emotion recognition project, using a fine-tuned
ResNet50 as the base model. All experiments were carried out on the AI-Lab SLURM
cluster using NVIDIA L4 GPUs, on the \texttt{feature/distributed-training} branch.

The central concepts explored are:

\paragraph{Amdahl's and Gustafson's Law.}
Amdahl's law bounds the maximum parallel speedup as
$S(p) = \frac{1}{(1-f) + f/p}$,
where $f$ is the parallelisable fraction and $p$ the number of processors.
For gradient-based deep learning, the serial fraction comes from data loading,
logging, and gradient synchronisation that cannot overlap with compute.
Gustafson's law provides a complementary view: as problem size grows with
more processors, the parallel fraction dominates and linear scaling becomes
achievable. Both laws together explain why small models like ResNet50 show
diminishing returns beyond a few GPUs, while large language models benefit
from thousands of GPUs.

\paragraph{Data Parallelism vs.\ Model Parallelism.}
Data parallelism replicates the full model on every GPU and splits the
input batch across them. After the backward pass, gradients are averaged
across all GPUs via AllReduce so every replica stays synchronised.
Communication cost scales with model size (one AllReduce per step).
Model parallelism instead splits the model itself across GPUs, which is
necessary when the model does not fit on a single device. It requires
activation transfers at every layer boundary during forward and backward
passes, making communication frequency much higher.
For ResNet50 ($\approx$100 MB, 25.6M parameters) on NVIDIA L4s (24 GB VRAM),
data parallelism is the natural choice: the model fits comfortably on one
GPU, so there is no pressure to split it.

\paragraph{Distributed Data Parallel (DDP).}
PyTorch's \texttt{DistributedDataParallel} wraps any model and handles
gradient synchronisation via NCCL (NVIDIA Collective Communications Library).
Each process is pinned to one GPU via \texttt{LOCAL\_RANK}, and a
\texttt{DistributedSampler} ensures each GPU sees a disjoint subset of the
training data each epoch. Only rank 0 logs metrics and saves checkpoints to
avoid duplicate writes.

\paragraph{Automatic Mixed Precision (AMP).}
AMP performs the forward and backward passes in FP16 while keeping the
master copy of weights in FP32. The smaller 16-bit tensors halve VRAM
usage, and on hardware with Tensor Cores the matrix multiplications can
execute faster. A \texttt{GradScaler} dynamically scales the loss to
prevent FP16 underflow during backpropagation.

\paragraph{ZeRO and DeepSpeed.}
The Zero Redundancy Optimizer (ZeRO) reduces memory pressure in multi-GPU
settings by partitioning different parts of the optimiser state across
GPUs rather than replicating it fully on each. Stage 1 partitions the
optimiser state (Adam moments and FP32 master weights). Stage 2 additionally
partitions the gradients. Stage 3 also partitions the model parameters
themselves, gathering them just-in-time for each forward pass.
ZeRO is designed for models with billions of parameters where the optimiser
state alone can be tens of gigabytes.

\paragraph{Federated Learning.}
Federated Learning trains a shared global model without centralising raw
data. Clients train locally on private data and send only model updates
(gradients or weights) to a server that aggregates them. This is
particularly relevant in privacy-sensitive domains such as medical imaging,
where regulatory constraints prevent data pooling. It was not implemented
in this project due to the single-institution academic setting, but
represents a natural extension for production deployments with
distributed data sources.

\paragraph{Neural Scaling Laws.}
Kaplan et al.\ (2020) showed that the loss of large neural language models
follows a power law in compute, dataset size, and parameter count:
$L(C) = (C_0 / C)^{\alpha}$.
The practical implication is that to halve the loss, one needs
$2^{1/\alpha}$ times more of the given resource—often an enormous
multiplier. These laws were derived for LLMs but provide a useful
order-of-magnitude frame of reference for any neural model.

% -----------------------------------------------------------------
\subsection{Implementation}

The distributed training implementation was developed in a dedicated
\texttt{feature/distributed-training} branch. The core script,
\texttt{train\_ddp.py}, extends the existing single-GPU \texttt{train.py}
with minimal changes to enable DDP across multiple GPUs on the same node
and across multiple nodes via \texttt{torchrun}.

\begin{lstlisting}[language=Python,
  caption={DDP initialisation in \texttt{train\_ddp.py}},
  label={lst:ddp_init},
  basicstyle=\ttfamily\small,
  keywordstyle=\color{blue},
  commentstyle=\color{gray},
  stringstyle=\color{red!60!black},
  breaklines=true]
import torch.distributed as dist
from torch.nn.parallel import DistributedDataParallel as DDP
from torch.utils.data import DistributedSampler

# Pin each process to its own GPU
torch.cuda.set_device(int(os.environ["LOCAL_RANK"]))
dist.init_process_group("nccl")

rank      = dist.get_rank()
local_rank = int(os.environ["LOCAL_RANK"])

# Wrap the model — gradient sync is automatic from here
model    = ResNet50FineTuned(model_config).to(local_rank)
ddp_model = DDP(model, device_ids=[local_rank])

# DistributedSampler gives each GPU a unique data shard
train_sampler = DistributedSampler(train_set)
train_loader  = DataLoader(
    train_set,
    batch_size=32,          # per-GPU batch size
    sampler=train_sampler,
    pin_memory=True,
    shuffle=False,
)

for epoch in range(epochs):
    train_sampler.set_epoch(epoch)   # re-shuffle each epoch
    for inputs, targets in train_loader:
        inputs, targets = inputs.to(local_rank), targets.to(local_rank)
        loss = criterion(ddp_model(inputs), targets)
        optimizer.zero_grad()
        loss.backward()          # AllReduce happens here
        optimizer.step()

    if rank == 0:                # only rank 0 saves checkpoints
        torch.save(model.state_dict(), "checkpoint.pt")
\end{lstlisting}

AMP was integrated into both single-GPU (\texttt{train.py}) and DDP training
using PyTorch's \texttt{torch.cuda.amp} module:

\begin{lstlisting}[language=Python,
  caption={AMP integration with GradScaler},
  label={lst:amp},
  basicstyle=\ttfamily\small,
  keywordstyle=\color{blue},
  commentstyle=\color{gray},
  stringstyle=\color{red!60!black},
  breaklines=true]
from torch.cuda.amp import autocast, GradScaler

scaler = GradScaler()

for inputs, targets in train_loader:
    with autocast():                   # FP16 forward + backward
        outputs = model(inputs)
        loss    = criterion(outputs, targets)

    scaler.scale(loss).backward()      # scaled gradients
    scaler.step(optimizer)
    scaler.update()
    optimizer.zero_grad()
\end{lstlisting}

The DeepSpeed ZeRO optimizer was configured via separate JSON configs for
stages 1, 2 and 3, run in FP32 to isolate the effect of partitioning from
any FP16-specific behaviour. The Jenkins pipeline on the
\texttt{feature/distributed-training} branch includes a
\texttt{Train Model (DDP)} stage gated by a parameter flag, so the
distributed path is opt-in without breaking the standard single-GPU
production flow.

% -----------------------------------------------------------------
\subsection{Experimental Setup}

All runs trained the same model (ResNet50FineTuned, 25.6M parameters) on
FER2013 ($\approx$50k images, 7 emotion classes, split 81/9/10 into
train/val/test) for 4 epochs. Hardware was NVIDIA L4 GPUs on the AI-Lab
SLURM cluster, with allocations from 1 GPU on one node up to 4 GPUs
across 2 nodes. Comparisons use \emph{steady-state per-epoch time}
(average of epochs 2–4) to exclude a one-time model weight download and
NCCL warm-up that inflated epoch 1 by roughly 1500 seconds on the
single-GPU baseline.

\begin{table}[h]
\centering
\caption{Full results across all eight experimental configurations.}
\label{tab:full_results}
\begin{tabular}{lcccccc}
\toprule
\textbf{Run} & \textbf{GPUs} & \textbf{AMP} & \textbf{ZeRO}
  & \textbf{Steady epoch (s)} & \textbf{Throughput (img/s)}
  & \textbf{VRAM (MB)} \\
\midrule
1-GPU FP32      & 1  & --  & -- & 170 &  76.2 & 2{,}906 \\
1-GPU AMP       & 1  & \checkmark & -- & 117 &  86.5 & 1{,}669 \\
DDP 2-GPU       & 2  & \checkmark & -- &  74 & 234.6 & 1{,}826 \\
DDP 4-GPU       & 4  & \checkmark & -- &  49 & 263.4 & 1{,}829 \\
DDP multi-node  & 4* & \checkmark & -- &  92 & 191.1 & 1{,}829 \\
ZeRO Stage 1    & 2  & --  & 1  & 223 & 178.7 & 4{,}871 \\
ZeRO Stage 2    & 2  & --  & 2  & 226 & 176.2 & 6{,}774 \\
ZeRO Stage 3    & 2  & --  & 3  & 170 & 122.5 & 4{,}962 \\
\bottomrule
\end{tabular}
\vspace{1mm}
{\small * 2 nodes $\times$ 2 GPUs each = 4 GPUs total.
All runs achieved 68–70\% validation accuracy.}
\end{table}

% -----------------------------------------------------------------
\subsection{D3.1 — Estimated Speedup by Parallelisation (Amdahl's Law)}

For ResNet50 training, the parallelisable fraction is estimated at
$f = 0.95$, covering the compute-bound forward and backward passes.
The serial fraction $s = 0.05$ covers data loading, logging, and gradient
synchronisation that does not overlap with compute.
Amdahl's law gives the theoretical maximum speedup as:

\[
S(p) = \frac{1}{(1 - f) + f/p} = \frac{1}{0.05 + 0.95/p}
\]

\begin{table}[h]
\centering
\caption{Theoretical vs.\ measured parallel speedup.}
\label{tab:amdahl}
\begin{tabular}{cccc}
\toprule
\textbf{GPUs ($p$)} & \textbf{Theoretical max ($S$)}
  & \textbf{Measured} & \textbf{Efficiency} \\
\midrule
2 & 1.82$\times$ & 1.58$\times$ & 87\% \\
4 & 3.33$\times$ & 2.39$\times$ & 72\% \\
8 & 5.93$\times$ & — & — \\
\bottomrule
\end{tabular}
\end{table}

Efficiency drops from 87\% at 2 GPUs to 72\% at 4 GPUs because the serial
fraction—gradient AllReduce, sampler synchronisation, and validation
pass—takes up a proportionally larger share of each iteration as compute
time shrinks. This is exactly what Amdahl's law predicts: there is a
diminishing return to adding more GPUs, and at 4 GPUs the overhead already
accounts for 28\% of wall time. A projected 8-GPU run would lose even more.

The measured efficiency gap also reflects the specifics of this setup: a
small model (25.6M params) with a small per-GPU batch (32 images) means
the gradient AllReduce, although operating on $\sim$100 MB of gradients,
represents a large fraction of the total step time. For larger models or
bigger batches, the compute-to-communication ratio improves and measured
efficiency would be closer to the theoretical bound.

% -----------------------------------------------------------------
\subsection{D3.2 — Scaling Needed to Halve the Test Loss}

Neural scaling laws (Kaplan et al.\ 2020) model the relationship between
loss and a resource $R$ as:
\[
L(R) = \left(\frac{R_0}{R}\right)^{\alpha}
\]
To halve the current loss, the required multiplier is $2^{1/\alpha}$.
Using the best observed validation loss of 0.83 as a baseline, and
applying the LLM-derived exponents as a rough frame of reference:

\begin{table}[h]
\centering
\caption{Resources required to halve the current validation loss (0.83)
  under the power-law scaling model.}
\label{tab:scaling_laws}
\begin{tabular}{lccc}
\toprule
\textbf{Resource} & \textbf{Exponent} & \textbf{Multiplier needed}
  & \textbf{Concrete scale} \\
\midrule
Compute (FLOPs) & $\alpha = 0.05$ & $\approx$1{,}000{,}000$\times$
  & $10^6\times$ current budget \\
Dataset size    & $\beta  = 0.095$ & $\approx$1{,}700$\times$
  & $\sim$68M images (vs.\ $\sim$40k) \\
Parameter count & $\gamma = 0.076$ & $\approx$8{,}200$\times$
  & $\sim$210B params (vs.\ 25.6M) \\
\bottomrule
\end{tabular}
\end{table}

These numbers are intentionally large. The correct interpretation is not
that the project should scale this far, but that ResNet50 fine-tuned on
FER2013 is already close to the compute-optimal frontier for this task.
Further loss reduction requires a qualitatively different approach—a
larger and more diverse dataset, stronger data augmentation, or a
different model architecture—rather than more of the same training
compute.

An important caveat: the scaling exponents ($\alpha$, $\beta$, $\gamma$)
were derived empirically for large language models trained on text.
Extrapolating them to a 25.6M-parameter vision model fine-tuned on 40k
images is rough. The estimates give a useful sense of order of magnitude
but should not be treated as precise budgets.

% -----------------------------------------------------------------
\subsection{D3.3 — Effect of Multi-GPU Data Parallelism}

Distributed Data Parallel training was implemented using
\texttt{torchrun} to launch one process per GPU. Each process holds a
full copy of the model and processes a distinct shard of each mini-batch.
After the backward pass, NCCL AllReduce averages gradients across all
ranks so every GPU takes an identical optimiser step.

The baseline for this comparison is the 1-GPU AMP configuration, since
all DDP runs also used AMP.

\begin{table}[h]
\centering
\caption{Multi-GPU DDP speedup (steady-state epoch time, epochs 2–4).}
\label{tab:ddp_speedup}
\begin{tabular}{lcccc}
\toprule
\textbf{Config} & \textbf{Steady epoch (s)}
  & \textbf{Throughput (img/s)} & \textbf{Speedup} & \textbf{Val acc} \\
\midrule
1-GPU AMP (reference) & 117 &  86.5 & 1.00$\times$ & 69.4\% \\
DDP 2-GPU             &  74 & 234.6 & 1.58$\times$ & 69.5\% \\
DDP 4-GPU             &  49 & 263.4 & 2.39$\times$ & 68.6\% \\
\bottomrule
\end{tabular}
\end{table}

Scaling is sub-linear but practically useful: 2 GPUs deliver a 1.58$\times$
speedup (87\% efficiency) and 4 GPUs a 2.39$\times$ speedup (72\% efficiency).
The diminishing returns from 2 to 4 GPUs are consistent with Amdahl's law
at $f = 0.95$. The primary bottleneck is the gradient AllReduce step,
which transfers $\sim$100 MB of gradients per training step.
For a small model and small per-GPU batch size (32), this communication
cost represents a larger share of total step time than it would for a
larger model, which explains the gap between theoretical and measured
efficiency.

Importantly, validation accuracy is stable across all configurations
(68.6–69.5\%), confirming that DDP does not degrade model quality.

% -----------------------------------------------------------------
\subsection{D3.4 — Effect of Multi-Node Parallelism}

The 4-GPU configuration was tested two ways: all GPUs on a single node
(1~node $\times$ 4 GPUs) and split across two nodes (2~nodes $\times$
2 GPUs). This isolates the effect of inter-node network communication on
training throughput.

\begin{table}[h]
\centering
\caption{Single-node vs.\ multi-node performance with 4 GPUs total.}
\label{tab:multinode}
\begin{tabular}{lcccc}
\toprule
\textbf{Config} & \textbf{GPUs} & \textbf{Steady epoch (s)}
  & \textbf{Throughput (img/s)} & \textbf{Penalty} \\
\midrule
DDP single-node & 4 (1$\times$4) & 49 & 263.4 & — \\
DDP multi-node  & 4 (2$\times$2) & 92 & 191.1 & 1.88$\times$ slower \\
\bottomrule
\end{tabular}
\end{table}

The multi-node configuration ran 1.88$\times$ slower than the equivalent
single-node setup, despite using the same total number of GPUs.
The cause is inter-node network bandwidth: the AI-Lab cluster uses standard
Ethernet between nodes (likely 10–25 GbE) rather than a high-bandwidth
fabric like InfiniBand, so the AllReduce step takes significantly longer
when gradients must traverse the network. Within a node, GPU-to-GPU
communication happens over PCIe (and potentially NVLink), which is an
order of magnitude faster.

For ResNet50, going multi-node is therefore counterproductive.
Multi-node parallelism only pays off when the model or its optimiser state
is too large to fit on a single node's GPUs, forcing cross-node
communication regardless. Since ResNet50 is $\sim$100 MB and a single
AI-Lab node has at least 24 GB of VRAM per GPU, a single-node solution
is always preferable here.

% -----------------------------------------------------------------
\subsection{D3.5 — Effect of Automatic Mixed Precision}

AMP was integrated into \texttt{train.py} using PyTorch's
\texttt{torch.cuda.amp.autocast} context manager and \texttt{GradScaler}.
The master copy of all parameters remains in FP32; only the forward and
backward pass tensors are cast to FP16. Dynamic loss scaling prevents
gradient underflow in half precision.

\begin{table}[h]
\centering
\caption{Effect of AMP on VRAM usage, training speed, and accuracy.}
\label{tab:amp}
\begin{tabular}{lcccc}
\toprule
\textbf{Precision} & \textbf{VRAM (MB)} & \textbf{Steady epoch (s)}
  & \textbf{Throughput (img/s)} & \textbf{Val accuracy} \\
\midrule
FP32 (baseline) & 2{,}906 & 170 & 76.2 & 70.0\% \\
AMP FP16        & 1{,}669 & 117 & 86.5 & 69.4\% \\
\midrule
\textbf{Change} & $-$43\% & $-$31\% & $+$14\% & $-$0.6\% \\
\bottomrule
\end{tabular}
\end{table}

AMP reduces VRAM consumption by 43\% (from 2906 MB to 1669 MB)—almost
halving the memory footprint without any change to the model architecture.
Epoch time drops by 31\% and throughput improves by 14\%.

The speed improvement is modest relative to the VRAM savings because
ResNet50 is not large enough to saturate the L4's Tensor Cores fully in
FP16. For larger models where Tensor Cores are the bottleneck, AMP can
deliver close to 2$\times$ speed improvement. The VRAM savings, however,
are independent of model size and are real regardless.

The accuracy delta of 0.6 percentage points is within run-to-run
noise and does not represent a meaningful quality loss. AMP is therefore
the single highest-value change that can be made to the training pipeline:
it requires only a \texttt{GradScaler} and \texttt{autocast}, no new
infrastructure, and frees up VRAM for larger batch sizes or more complex
models.

% -----------------------------------------------------------------
\subsection{D3.6 — Effect of ZeRO Optimizer Stages}

DeepSpeed ZeRO was evaluated with stages 1, 2 and 3 on 2 GPUs, using FP32
throughout to isolate the effect of state partitioning from precision effects.
The reference baseline is the 2-GPU DDP+AMP configuration.

\begin{table}[h]
\centering
\caption{ZeRO optimizer stages: training speed and VRAM usage.}
\label{tab:zero}
\begin{tabular}{lcccc}
\toprule
\textbf{Stage} & \textbf{Steady epoch (s)}
  & \textbf{Throughput (img/s)} & \textbf{VRAM peak (MB)} & \textbf{Notes} \\
\midrule
No ZeRO (DDP+AMP) &  74 & 234.6 & 1{,}826 & Reference \\
ZeRO Stage 1 (FP32) & 223 & 178.7 & 4{,}871 & Higher VRAM than ref. \\
ZeRO Stage 2 (FP32) & 226 & 176.2 & 6{,}774 & Highest of all configs \\
ZeRO Stage 3 (FP32) & 170 & 122.5 & 4{,}962 & Slowest throughput \\
\bottomrule
\end{tabular}
\end{table}

The results appear counterintuitive—ZeRO is designed to \emph{reduce}
VRAM, yet all three stages consume \emph{more} memory than the DDP+AMP
baseline. The explanation lies in the mismatch between model scale and
framework overhead.

ZeRO is engineered for models with billions of parameters, where the Adam
optimiser state (first and second moments plus FP32 master weights) can
reach tens of gigabytes. For ResNet50, the total optimiser state is only
$\sim$200 MB. Partitioning it across 2 GPUs saves approximately 100 MB
per GPU. Meanwhile, DeepSpeed pre-allocates fixed-size communication
buffers that do not scale down with the model:
\texttt{reduce\_bucket\_size} ($\sim$2 GB) for Stage 2's ReduceScatter,
and \texttt{stage3\_prefetch\_bucket\_size} ($\sim$2 GB) for Stage 3's
AllGather. These buffers far exceed the partitioning savings, resulting
in a net increase in VRAM rather than a decrease.

Stage 3 is also the slowest configuration (122.5 img/s vs.\ 234.6 img/s
for DDP+AMP), because it must AllGather all model parameters before every
forward pass and re-partition them afterwards. For a model that fits easily
on one GPU, this is pure overhead.

ZeRO becomes genuinely useful at the 1B+ parameter scale, where the model
and optimiser state cannot fit on a single GPU and the fixed buffer
overhead is amortised over savings that are orders of magnitude larger.
For this project, the ZeRO implementation is retained in the
\texttt{feature/distributed-training} branch as a reference and
future-proofing measure, but is not promoted to the production pipeline.

% -----------------------------------------------------------------
\subsection{Integration into the MLOps Pipeline}

AMP has been integrated directly into the production training script
(\texttt{train.py}) via a configuration flag (\texttt{train.amp: true}
in \texttt{config/train\_config.yaml}). The Jenkins pipeline continues
to invoke \texttt{python3 train.py} with a single GPU, preserving the
existing workflow while benefiting from the 43\% VRAM reduction and 14\%
throughput improvement.

The DDP training path (\texttt{train\_ddp.py}) is available as an
opt-in stage in the Jenkinsfile, gated behind a build parameter. When
activated, it replaces the \texttt{Train Model} stage with a
\texttt{torchrun}-based launch across the requested number of GPUs.
This design keeps the standard pipeline fast and simple while making
multi-GPU training available for larger experiments.

% -----------------------------------------------------------------
\subsection{Summary}

\begin{table}[h]
\centering
\caption{Summary of all D3 findings.}
\label{tab:d3_summary}
\begin{tabular}{lp{9cm}}
\toprule
\textbf{Item} & \textbf{Finding} \\
\midrule
D3.1 & DDP achieves 1.58$\times$ at 2 GPUs (87\% efficiency) and 2.39$\times$
       at 4 GPUs (72\% efficiency) relative to a single-GPU AMP baseline,
       consistent with Amdahl's law at $f = 0.95$. \\
D3.2 & Halving the current validation loss (0.83) would require
       $\sim$1{,}700$\times$ more training data ($\sim$68M images),
       $\sim$8{,}200$\times$ more parameters ($\sim$210B), or
       $\sim$10^6$\times$ more compute, indicating diminishing returns
       from incremental scaling of the current approach. \\
D3.3 & 2-GPU DDP reduces steady-state epoch time from 117 s to 74 s
       (1.58$\times$); 4-GPU DDP reduces it further to 49 s (2.39$\times$).
       Validation accuracy is stable at 68–70\% across all configurations. \\
D3.4 & Spreading 4 GPUs across 2 nodes is 1.88$\times$ slower than the
       same 4 GPUs on one node (92 s/epoch vs.\ 49 s/epoch),
       due to the slower inter-node Ethernet fabric on AI-Lab. \\
D3.5 & AMP reduces VRAM by 43\% (2906 MB $\to$ 1669 MB) and speeds up
       training by 14\% with no meaningful accuracy cost ($-$0.6\%).
       It is the single most impactful change for this model size. \\
D3.6 & ZeRO Stages 1–3 all increased VRAM relative to the DDP+AMP
       baseline because DeepSpeed's fixed communication buffers
       ($\sim$2 GB each) vastly exceed the savings from partitioning a
       200 MB optimiser state. ZeRO is counterproductive at this scale
       and is kept off the production pipeline. \\
\bottomrule
\end{tabular}
\end{table}
